% ============ 08. Género y diversidad ============
\section{Género y diversidad en la ciencia}

\subsection{Qué dice la literatura}

El referente clásico es el estudio global de \citet{lariviere2013global}
(\emph{Nature}), que documentó la subrepresentación femenina y su variación por
disciplina. La literatura reciente en español y portugués aporta evidencia
regional: análisis bibliométricos comparativos de la \textbf{Comunidad Andina}
(2020--2024) sobre disparidades de género en producción (sección 13.2) y
estudios sobre participación femenina en STEM (sección 13.2). El enfoque
\textbf{interseccional} \citep{crenshaw1989demarginalizing} recomienda cruzar
género con etnia, territorio y otras dimensiones, en vez de analizar el género
de forma aislada. Las métricas estándar son proporciones de participación por
disciplina, brechas relativas (odds ratios) y razones de representación frente
a la población de referencia; la literatura advierte contra comparar
proporciones sin ajustar por composición de área y cohorte.

\subsection{Relación con este repositorio}

El proyecto mide la participación femenina por gran área OCDE (hallazgo: 24\,\%
en Ingeniería vs.\ 48\,\% en ciencias médicas), la evolución temporal y la
brecha en productividad (5 de 6 grandes áreas con brecha); en diversidad,
compara grupos étnicos, discapacidad y víctimas frente al censo DANE 2018 y el
RUV (afro 3$\times$, indígena 7.9$\times$, discapacidad 8$\times$
subrepresentados) y produce un cruce interseccional género$\times$etnia. La
auditoría verificó debilidades: variables auto-declaradas y capturadas solo
desde 2021 (sin serie temporal); comparaciones sin modelos ajustados (odds
ratios con IC); el bug metodológico de \texttt{comparar\_dane()}: el comentario
promete excluir "NINGUN GRUPO ETNICO" del denominador pero el código solo
excluye "NO DISPONIBLE", con lo que la subrepresentación queda diluida (el
arreglo recomendado es excluir de verdad ese grupo del denominador); y el cruce
interseccional frágil ante subgrupos sin mujeres.

\subsection{Mejoras recomendadas}

\begin{enumerate}
  \item Formalizar la inferencia de brechas con \textbf{modelos logísticos
        ajustados por área, cohorte y categoría} (odds ratios con IC), como
        recomienda la literatura \citep{lariviere2013global} y sus extensiones
        regionales (sección 13.2).
  \item Reportar \textbf{razones de representación con intervalos de
        confianza} (bootstrap) y declarar explícitamente la ventana 2021 como
        única comparable para diversidad.
  \item Robustecer el \textbf{análisis interseccional} (género$\times$etnia
        $\times$área) con \texttt{reindex} de columnas y tamaños de celda
        mínimos (evitar desagregaciones sin respaldo muestral), siguiendo el
        enfoque de \citet{crenshaw1989demarginalizing}.
\end{enumerate}
